\documentclass[11pt]{article}

\usepackage[T1]{fontenc}
\usepackage{lmodern}
\usepackage{microtype}
\usepackage[margin=1in]{geometry}
\usepackage{amsmath,amssymb,amsthm,mathtools}
\usepackage{booktabs}
\usepackage{enumitem}
\usepackage{xcolor}
\usepackage[colorlinks=true,linkcolor=blue!55!black,citecolor=green!40!black,urlcolor=blue!55!black]{hyperref}

\setlength{\emergencystretch}{3em}
\setlist{nosep,leftmargin=2em}
\allowdisplaybreaks

\newtheorem{theorem}{Theorem}[section]
\newtheorem{lemma}[theorem]{Lemma}
\newtheorem{corollary}[theorem]{Corollary}
\newtheorem{proposition}[theorem]{Proposition}
\theoremstyle{definition}
\newtheorem{definition}[theorem]{Definition}
\theoremstyle{remark}
\newtheorem{remark}[theorem]{Remark}

\newcommand{\dist}{\operatorname{dist}}
\newcommand{\diam}{\operatorname{diam}}
\newcommand{\Pot}{\operatorname{Pot}}
\newcommand{\Imp}{\operatorname{Imp}}
\newcommand{\IR}{\operatorname{IR}}
\newcommand{\poly}{\operatorname{poly}}
\newcommand{\eps}{\varepsilon}
\newcommand{\N}{N}
\newcommand{\LOC}{\textnormal{\textsc{Loc}}}
\newcommand{\LOP}{\textnormal{\textsc{Lop}}}
\newcommand{\LOPE}{\textnormal{\textsc{Lope}}}
\newcommand{\LOCAL}{\textnormal{\textsc{Local}}}

\title{Direct Sequence Localization for Distributed Potential Problems}
\author{}
\date{}

\begin{document}
\maketitle
\vspace{-2.5em}

\begin{abstract}
We study distributed local search in the deterministic \LOCAL\ model.  Balliu,
Boudier, d'Amore, Kuhn, Olivetti, Schmid, and Suomela introduced locally optimal
problems and showed that every bounded-degree problem governed by a local potential
admits a polylogarithmic-round algorithm.  Their analysis localizes each improving
move separately and then localizes a sequence of moves in an auxiliary interaction
graph.

The main observation here is that the whole sequence can instead be replayed inside
one ball of the original graph.  Minimality pays for the discarded part of every
move, and the only remaining error is charged to potential terms crossing the ball
boundary.  For a radius-$\rho$ potential, normalized local scale $K$, and normalized
upper bound $P$ on the global potential, this gives a deterministic algorithm in
\[
 O_{\rho}\!\left(
 \left[1+\min\left\{
 K(\Delta^\rho+1)\log n\log(2P),
 \sqrt{Kn\log n\log(2P)}
 \right\}\right]
 \log^2 n\,\poly(\log\log n)
 \right)
\]
rounds.  Thus, for fixed local scale, the polynomial part is
$\min\{\Delta^\rho,\sqrt n\}$.

For edge potentials the boundary is sharper.  In particular, ordinary locally
optimal cut has a deterministic
\[
 O\!\left(\min\{\Delta\log^4 n,\sqrt n\log^3 n\}
 \poly(\log\log n)\right)
\]
algorithm.  Together with the randomized and quantum lower bound of Balliu et al.,
this determines the polynomial dependence on $\Delta$ and $n$ up to logarithmic
factors.  As supporting evidence that the boundary-volume dependence of such generic
local-search algorithms is inherent, we also analyze the stronger task in which no
subset of a radius-$s$ ball may improve a cut.  For every fixed $s\geq 1$, its
deterministic upper bound and randomized/quantum lower bound have the same polynomial
part
\[
 \min\{\Delta^{s+1},\sqrt{\Delta n},n^{2/3}\}.
\]
\end{abstract}

\section{Introduction}

In an undirected graph $G=(V,E)$, a cut is a labeling $x:V\to\{0,1\}$.
It is \emph{locally optimal} if flipping any one vertex does not increase the
number of cut edges.  Equivalently, every vertex has at least as many neighbors on
the other side of the cut as on its own side.  A centralized improvement process is
immediate: repeatedly flip an improving vertex.  Since each flip increases the cut
by at least one edge, the process stops after at most $|E|$ flips.  The challenge is
to parallelize this process without creating a long cascade of mutually dependent
repairs.

We work in the standard synchronous \LOCAL\ model
\cite{Linial1992,NaorStockmeyer1995}: vertices have distinct identifiers, messages
have unbounded size, and after $T$ rounds a vertex can know its entire radius-$T$
neighborhood.  Local computation is free.  The problem of finding a locally optimal
cut will be abbreviated by $\LOC$.

Balliu et al.~\cite{BalliuEtAl2026} introduced a broad framework of locally optimal
problems (\LOP s).  A global potential is written as a sum of bounded-radius local
terms, and an invalid vertex has a one-block relabeling that strictly decreases this
potential.  They obtained the first polylogarithmic-round algorithms for all
bounded-degree \LOP s.  For $\LOC$, their deterministic complexity is
$O(\Delta^2\log^8 n\,\poly(\log\log n))$, and they proved an
$\Omega(\min\{\Delta,\sqrt n\})$ lower bound for randomized and quantum-\LOCAL
algorithms in the usual range $\Delta\leq cn$, for any fixed $c<1$.

The present work improves the upper bound by changing the localization argument.
The prior proof first replaces a large improving move by a small-diameter move and
then represents moves as vertices of an auxiliary interaction graph.  The two
localization steps multiply their losses.  Here the entire sequence is replayed
directly inside a ball of $G$.  The proof counts occurrences of original vertices,
so the auxiliary graph is unnecessary.

\subsection{What is new and what is inherited}

The definitions of \LOP, improving sets, improving ratios, and minimal improving
sets, as well as the threshold-scheduling paradigm, originate in
Balliu et al.~\cite{BalliuEtAl2026}.  We also use the deterministic strong-diameter
network decomposition of Ghaffari and Grunau~\cite{GhaffariGrunau2024}; this line of
work began with the polylogarithmic deterministic decomposition of Rozho\v{n} and
Ghaffari~\cite{RozhonGhaffari2020}.

The new ingredients and conclusions are:
\begin{enumerate}
  \item a direct replay lemma for a whole sequence of minimal improving moves;
  \item a one-step sequence-localization theorem in the original graph;
  \item a deterministic generic \LOP\ algorithm whose polynomial dependence improves
  from $\Delta^{2\rho}$ to $\min\{\Delta^\rho,\sqrt n\}$ at fixed local scale;
  \item an edge-boundary specialization giving a deterministic
  $\widetilde O(\min\{\Delta,\sqrt n\})$ algorithm for $\LOC$, which matches the
  known lower bound up to logarithmic factors; and
  \item a supporting upper/lower-bound pair for radius-$s$ local cut.  The latter is
  not needed for the $\LOC$ algorithm; its role is to show that the polynomial
  boundary-volume scales encountered by generic local-search methods occur on
  concrete simple unweighted graphs.
\end{enumerate}

The distinction is summarized below.  The table suppresses logarithmic factors and
fixed-radius constants.
\begin{center}
\begin{tabular}{@{}lll@{}}
\toprule
Problem & Previously known & Result here \\
\midrule
General radius-$\rho$ \LOP
  & $\widetilde O(K^2\Delta^{2\rho})$ deterministic
  & $\widetilde O_\rho(\min\{K\Delta^\rho,\sqrt{Kn}\})$ deterministic \\
Ordinary $\LOC$
  & $\widetilde O(\Delta^2)$ deterministic
  & $\widetilde O(\min\{\Delta,\sqrt n\})$ deterministic \\
Radius-$s$ local cut
  & --
  & $\widetilde\Theta_s(\min\{\Delta^{s+1},\sqrt{\Delta n},n^{2/3}\})$ \\
\bottomrule
\end{tabular}
\end{center}
The $\widetilde\Theta_s$ notation in the last row compares a deterministic upper
bound with a randomized and quantum lower bound, and hence refers only to the
polynomial dependence on $n$ and $\Delta$.

\section{Model and normalized potential parameters}

For $A\subseteq V$, let $\N_t[A]=\{u:\dist(u,A)\leq t\}$, and write
$\N_t[v]=\N_t[\{v\}]$.  A radius-$\rho$ local-potential problem has a finite output
alphabet and a potential
\begin{equation}\label{eq:potential}
  \Pot_G(\ell)=\sum_{z\in V}\phi_z(\ell|_{\N_\rho[z]}),
  \qquad 0\leq \phi_z\leq \Lambda.
\end{equation}
The input labels and the topology inside $\N_\rho[z]$ may be arguments of $\phi_z$;
we suppress them.  A ``vertex block'' may contain the output of a vertex and its
incident half-edge outputs.  Thus this notation includes the formal \LOP\ model of
\cite{BalliuEtAl2026}; all that is used below is that legal block relabelings are
finite and restrictions of a legal multi-block relabeling are legal.

We assume that whenever the current output is invalid at a vertex, some relabeling
of its block decreases \eqref{eq:potential} by at least $\lambda>0$.  Divide all
potentials by $\lambda$ and define
\begin{equation}\label{eq:KP}
  K=\frac{\Lambda}{\lambda},
  \qquad
  P\geq \max_{\ell}\frac{\Pot_G(\ell)}{\lambda}.
\end{equation}
Henceforth potentials are normalized so that the guaranteed correction is at least
one, each local term lies in $[0,K]$, and every global potential is at most $P$.
The algorithm only needs a common upper bound on $P$.  The universal choice $P=Kn$
is available for \eqref{eq:potential}; retaining $P$ makes the dependence on the
maximum global potential explicit.  The factor $K$ is the one unavoidable local
scale parameter.  For a fixed problem family at bounded degree, it is a constant.

\begin{definition}[Improving move]
For a labeling $\ell$, a move $M=(A,\ell_A)$ consists of a nonempty support
$A\subseteq V$ and a legal relabeling supported on $A$.  If $\ell^M$ is the
resulting labeling, define
\[
  \Imp_\ell(M)=\Pot_G(\ell)-\Pot_G(\ell^M),
  \qquad
  \IR_\ell(M)=\frac{\Imp_\ell(M)}{|A|}.
\]
The move is improving if $\Imp_\ell(M)>0$.  An improving move of ratio $q$ is
\emph{minimal} if every move supported on a proper nonempty subset of $A$ has ratio
at most $q$.
\end{definition}

Every improving move contains a minimal improving move of at least the same ratio:
choose, among all improving moves supported inside its support, one of maximum ratio
and then inclusion-minimal support.  A sequence $M_1,\ldots,M_h$ is
$\beta$-improving if $M_i$ is evaluated after the preceding moves and has ratio at
least $\beta$.

Let
\begin{equation}\label{eq:ballsize}
  B_\rho(\Delta)=1+\Delta\sum_{j=0}^{\rho-1}(\Delta-1)^j
  =O_\rho(\Delta^\rho)
\end{equation}
be the maximum size of a radius-$\rho$ ball in a graph of maximum degree $\Delta$.

\section{Direct localization of an improving sequence}

Fix an initial labeling $\ell_0$ and a sequence of minimal $\beta$-improving moves
$M_i=(A_i,\ell_i)$, where $\ell_i$ is the labeling after move $i$.  Count repeated
changes by
\[
  w(u)=|\{i:u\in A_i\}|,
  \qquad w(U)=\sum_{u\in U}w(u).
\]
The sequence decreases the normalized potential by at least $\beta w(V)$, while the
initial potential is at most $P$.  Therefore
\begin{equation}\label{eq:totalmass}
  w(V)\leq \frac{P}{\beta}.
\end{equation}

For $U\subseteq V$, the \emph{replay inside $U$} starts from $\ell_0$, follows the
labels of the true sequence on $U$, and permanently freezes all blocks outside $U$
at their labels in $\ell_0$.  Thus replay step $i$ applies only $A_i\cap U$.

\begin{lemma}[Direct replay]\label{lem:replay}
Let $\ell'_h$ be the final labeling obtained by replaying a minimal
$\beta$-improving sequence inside $U$.  Then
\[
 \Pot_G(\ell_0)-\Pot_G(\ell'_h)
 \geq
 \beta w(U)-2K B_\rho(\Delta)
 \sum_{\substack{u\in U\\ \dist(u,V\setminus U)\leq 2\rho}}w(u).
\]
More precisely, every occurrence of a boundary vertex can be charged only for local
terms centered in $\N_\rho[u]$, and each such term contributes discrepancy at most
$2K$.
\end{lemma}

\begin{proof}
Consider step $i$ and put $B=A_i\cap U$ and $C=A_i\setminus U$.  Let $q_i$ be the
ratio of the full move.  In the true state immediately before step $i$, first apply
the prescribed part on $C$.  Minimality says that its improvement is at most
$q_i|C|$; this is also true if the improvement is nonpositive.  Since the full move
improves by $q_i(|B|+|C|)$, the subsequent prescribed relabeling on $B$ has marginal
improvement at least $q_i|B|$.

Compare this marginal with replay step $i$.  Before and after the step, the two
worlds agree on $U$.  A local term can have different marginals only if its
radius-$\rho$ scope contains both a changed vertex of $B$ and a vertex outside $U$.
Such a changed vertex lies within distance $2\rho$ of $V\setminus U$.  A fixed
occurrence of $u$ belongs to at most $B_\rho(\Delta)$ local scopes, and the difference
between two marginals of one term in $[0,K]$ is at most $2K$.  Hence replay step $i$
improves by at least
\[
 q_i|B|-2K B_\rho(\Delta)
 |\{u\in B:\dist(u,V\setminus U)\leq2\rho\}|.
\]
Sum over $i$ and telescope the replay potential.  Since $q_i\geq\beta$, the first
sum is at least $\beta w(U)$, and the second sum is exactly the asserted occurrence
count.
\end{proof}

The proof intentionally overcounts a local term when several changed vertices affect
it.  Counting possible centers as well as changed occurrences gives a square-root
alternative.

\begin{theorem}[Sequence localization]\label{thm:localization}
Let $0<\eps<\beta$, and suppose a vertex $v$ occurs in a sequence of minimal
$\beta$-improving moves.  In the initial labeling there is a minimal improving move
of ratio at least $\beta-\eps$, contained in the union of the sequence and in a ball
around $v$ of radius
\begin{equation}\label{eq:localradius}
 L=O_\rho\!\left(
 1+\min\left\{
 \frac{K(\Delta^\rho+1)}{\eps}\log\!\left(1+\frac{P}{\beta}\right),
 \sqrt{\frac{Kn}{\eps}\log\!\left(1+\frac{P}{\beta}\right)}
 \right\}\right).
\end{equation}
No diameter assumption on the moves in the sequence is required.
\end{theorem}

\begin{proof}
Let $U_i=\N_{2\rho i}[v]$, $W_i=w(U_i)$, and for $i\geq1$ put
\[
 a_i=W_i-W_{i-1},
 \qquad
 Z_i=|U_{i+1}\setminus U_{i-1}|.
\]
Because $v$ occurs, $W_0\geq1$.  For every $k$,
\begin{equation}\label{eq:harmonic}
 \sum_{i=1}^k\frac{a_i}{W_i}
 \leq \sum_{i=1}^k\log\frac{W_i}{W_{i-1}}
 =\log\frac{W_k}{W_0}
 \leq\log\!\left(1+\frac{P}{\beta}\right),
\end{equation}
where the last inequality uses \eqref{eq:totalmass}.  The annuli
$U_{i+1}\setminus U_{i-1}$ overlap at most twice, so
\begin{equation}\label{eq:annulus-volume}
 \sum_{i=1}^k Z_i\leq2n.
\end{equation}

Replay inside $U_i$.  Any changed vertex within distance $2\rho$ of the complement
belongs to $U_i\setminus U_{i-1}$, whose total occurrence mass is $a_i$.  The center
of every discrepant local term lies in $U_{i+1}\setminus U_{i-1}$.  Repeating the
proof of Lemma~\ref{lem:replay}, but using the smaller of the changed-vertex and
center counts, shows that the replay loses at most
\begin{equation}\label{eq:two-boundaries}
 2K a_i\min\{B_\rho(\Delta),Z_i\}.
\end{equation}

Suppose the quantity in \eqref{eq:two-boundaries} exceeds $\eps W_i$ for every
$1\leq i\leq k$.  Using the first term in the minimum and
\eqref{eq:harmonic} gives
\[
 k=O_\rho\!\left(
 \frac{K(\Delta^\rho+1)}{\eps}
 \log\!\left(1+\frac{P}{\beta}\right)\right).
\]
Using the second term, then Cauchy--Schwarz and
\eqref{eq:annulus-volume}, gives
\[
 \log\!\left(1+\frac{P}{\beta}\right)
 \geq \frac{\eps}{2K}\sum_{i=1}^k\frac1{Z_i}
 \geq \frac{\eps k^2}{4Kn}.
\]
If some $Z_i=0$, the boundary error is already zero.  Consequently, within the
number of annuli displayed in \eqref{eq:localradius}, some $i$ has boundary loss at
most $\eps W_i$.

For this $i$, replay improves the initial potential by at least
$(\beta-\eps)W_i>0$.  Let $A$ be the set of blocks on which the final replay differs
from $\ell_0$.  Every vertex of $A$ occurs at least once, hence $|A|\leq W_i$; the
replay move has ratio at least $\beta-\eps$.  Taking a minimal improving submove
cannot decrease the ratio.  Its support is contained in $U_i$ and in the union of
the original sequence.
\end{proof}

\section{The generic deterministic algorithm}

We use a strong-diameter network decomposition: a coloring of clusters such that
adjacent clusters have different colors and every cluster has bounded diameter in
its induced subgraph.  Ghaffari and Grunau~\cite{GhaffariGrunau2024} give, in
$\widetilde O(\log^2 n)$ deterministic rounds, a decomposition with
$Q=O(\log n)$ colors and strong diameter $D=O(\log n)$.

Set
\[
 \eta=\frac1{8Q},
 \qquad
 R_j=\frac14+\frac{j-1}{4Q},
 \qquad 1\leq j\leq Q+1.
\]
Then $R_{j+1}-\eta>R_j$ and $R_{Q+1}=1/2$.  Choose $L$ from
Theorem~\ref{thm:localization} uniformly for $\beta\geq1/4$ and $\eps=\eta$, and let
$p=2(L+\rho)+1$.

Compute a network decomposition of the power graph $G^p$.  For a cluster $C$, let
its write region be $\N_L[C]$ and its observation region be $\N_{L+2\rho}[C]$.
Distinct clusters of one color have write regions at distance greater than $2\rho$,
so no local potential term is affected by two of them.  A cluster has diameter at
most $pD$ when power-graph edges are replaced by paths in $G$; its observation region
can therefore be gathered at a leader in $O(pD+L+\rho)$ rounds.

Process the $Q$ colors in order.  In a color-$j$ cluster, the leader repeatedly
chooses a maximum-ratio improving relabeling supported in the write region and
applies it while its ratio is at least $R_j$.  The leader can brute-force this finite
optimization because local computation is unbounded.  A maximum-ratio choice is
minimal.  The simulation terminates because the label space is finite and every
accepted move strictly decreases the potential.  Same-color simulations are
independent and may be serialized for analysis.

\begin{theorem}[Generic deterministic upper bound]\label{thm:generic}
Every normalized radius-$\rho$ local-potential problem described in
Section~2 can be solved deterministically in
\begin{equation}\label{eq:generic-bound}
 O_\rho\!\left(
 \left[1+\min\left\{
 K(\Delta^\rho+1)\log n\log(2P),
 \sqrt{Kn\log n\log(2P)}
 \right\}\right]
 \log^2 n\,\poly(\log\log n)
 \right)
\end{equation}
\LOCAL\ rounds.  If $K=O(1)$ and $P\leq n^{O(1)}$, this is
$\widetilde O_\rho(\min\{\Delta^\rho,\sqrt n\})$.
\end{theorem}

\begin{proof}
It remains to prove correctness and count rounds.  Suppose the final labeling is
invalid at $v$, and let $C$ be the decomposition cluster containing $v$, of color
$j$.  Consider the state just after color $j$ was processed.  Serialize every move
made in later color phases and append, in the final state, a correcting singleton
move at $v$.  Every actual later move is minimal and has ratio at least $R_{j+1}$;
the appended correction has ratio at least one and is minimal.  Thus these moves
form a minimal $R_{j+1}$-improving sequence in which $v$ occurs.

Apply Theorem~\ref{thm:localization} to this sequence.  In the state immediately
after color $j$, it produces a move of ratio at least
$R_{j+1}-\eta>R_j$ supported in $\N_L[v]\subseteq\N_L[C]$.  This contradicts
exhaustion of the write region during phase $j$.  Hence the final labeling is valid.

One round on $G^p$ costs at most $p$ rounds on $G$.  The decomposition costs
$p\widetilde O(\log^2 n)$ rounds.  The $Q$ color phases cost
$O(Q(pD+L+\rho))=O(p\log^2 n)$ rounds.  Substituting $\eps=\Theta(1/\log n)$ and
$\beta\geq1/4$ in \eqref{eq:localradius} yields \eqref{eq:generic-bound}.
\end{proof}

\begin{remark}[Comparison with prior work]
The prior deterministic bound of Balliu et al.~\cite{BalliuEtAl2026} has polynomial
part $K^2\Delta^{2\rho}$.  Theorem~\ref{thm:generic} removes one localization step,
one power of the local scale, and one power of the radius-$\rho$ ball volume.  It also
exposes the alternative $\sqrt{Kn}$ scale.  No matching lower bound is asserted for
every abstract \LOP; the tight statement below is for the concrete edge-potential
problem $\LOC$.
\end{remark}

\section{Edge potentials and a tight algorithm for locally optimal cut}

Suppose the potential is a sum over edges,
\[
 \Pot_G(\ell)=\sum_{e\in E}\psi_e(\ell|_e),
 \qquad 0\leq\psi_e\leq\Gamma.
\]
After division by the minimum correction $\lambda$, put $K_E=\Gamma/\lambda$ and
again let $P$ bound the normalized global potential.  Such problems are called
locally optimal problems on edges (\LOPE s) in~\cite{BalliuEtAl2026}.

\begin{lemma}[Edge replay]\label{lem:edge-replay}
For a sequence replayed inside $U$,
\[
 \Pot_G(\ell_0)-\Pot_G(\ell'_h)
 \geq \beta w(U)-2K_E\sum_{u\in U}w(u)\deg_{V\setminus U}(u).
\]
\end{lemma}

\begin{proof}
Repeat the proof of Lemma~\ref{lem:replay}.  A marginal can differ between the true
and frozen-outside executions only on an edge with one changed endpoint in $U$ and
the other endpoint outside $U$.  The discrepancy per crossing edge is at most
$2K_E$.  Charging it to the changed endpoint proves the formula.
\end{proof}

Take $U_i=\N_i[v]$, let $a_i$ be the occurrence mass of the $i$th layer, and put
$Z_i=|U_{i+1}\setminus U_{i-1}|$.  A changed boundary vertex has at most
$\min\{\Delta,Z_i\}$ neighbors outside $U_i$.  The proof of
Theorem~\ref{thm:localization} therefore gives an edge-potential localization radius
\begin{equation}\label{eq:edge-localization}
 O\!\left(1+\min\left\{
 \frac{K_E\Delta}{\eps}\log\!\left(1+\frac{P}{\beta}\right),
 \sqrt{\frac{K_E n}{\eps}\log\!\left(1+\frac{P}{\beta}\right)}
 \right\}\right).
\end{equation}
Using this radius in the scheduler from Section~4 proves the following edge version
of Theorem~\ref{thm:generic}.

\begin{theorem}[Edge-potential upper bound]\label{thm:edge-upper}
A normalized edge-potential problem can be solved deterministically in
\[
 O\!\left(
 \left[1+\min\left\{
 K_E\Delta\log n\log(2P),
 \sqrt{K_E n\log n\log(2P)}
 \right\}\right]
 \log^2 n\,\poly(\log\log n)
 \right)
\]
rounds.
\end{theorem}

For $\LOC$, use the potential
\[
 \Pot_G(x)=|\{uv\in E:x(u)=x(v)\}|,
\]
the number of uncut edges.  If $v$ is invalid, flipping $v$ decreases the potential
by a positive integer, hence by at least one.  Thus $K_E=1$, and the maximum global
potential is $P=|E|\leq n^2/2$.

\begin{corollary}[Tight deterministic algorithm for $\LOC$]\label{cor:loc}
Locally optimal cut can be solved deterministically in
\[
 O\!\left(
 \min\{\Delta\log^4 n,\sqrt n\log^3 n\}
 \poly(\log\log n)
 \right)
\]
rounds.
\end{corollary}

Balliu et al.~\cite[Theorem~5.2 and Corollary~5.6]{BalliuEtAl2026} proved that, for
every fixed $c<1$ and $\Delta\leq cn$, randomized and quantum-\LOCAL
algorithms with constant global success probability require
$\Omega(\min\{\Delta,\sqrt n\})$ rounds.  Corollary~\ref{cor:loc} therefore settles
the polynomial dependence on both parameters.  This tightness claim does not include
the remaining logarithmic factors.

\section{A supporting barrier: radius-local cuts}

This section is logically independent of the generic algorithm.  It studies a
stronger cut-stability notion in order to test whether the boundary-volume scales
behind direct localization can be improved for all local-search tasks.

\begin{definition}[Radius-$s$ locally optimal cut]
For a fixed integer $s\geq1$, a cut $x:V\to\{0,1\}$ is \emph{radius-$s$ locally
optimal} if, for every $v\in V$ and every $S\subseteq\N_s[v]$, simultaneously
flipping all vertices of $S$ does not increase the number of cut edges.
\end{definition}

Ordinary $\LOC$ corresponds to allowing only singleton flips, while radius-$s$
optimality quantifies over every coalition contained in one radius-$s$ ball.

\begin{theorem}[Radius-local cut bounds]\label{thm:radius-cut}
For every fixed $s\geq1$, radius-$s$ locally optimal cut has a deterministic
algorithm with complexity
\begin{equation}\label{eq:radius-upper}
 O_s\!\left(
 \min\{\Delta^{s+1}\log^4 n,
        \sqrt{\Delta n}\log^3 n,
        n^{2/3}\log^{8/3}n\}
 \poly(\log\log n)
 \right).
\end{equation}
Conversely, for all sufficiently large feasible parameter values, every randomized
or quantum-\LOCAL\ algorithm with global success probability at least $2/3$
requires
\begin{equation}\label{eq:radius-lower}
 \Omega_s\!\left(
 \min\{\Delta^{s+1},\sqrt{\Delta n},n^{2/3}\}
 \right)
\end{equation}
rounds, even on connected simple unweighted graphs.
\end{theorem}

The upper bound is included for completeness, but the main message for generic
algorithms is the lower bound.  In the degree-dominated regime
$\Delta\leq n^{1/(2s+1)}$, the $\Omega_s(\Delta^{s+1})$ term exactly matches the
volume scale of an edge interaction combined with a radius-$s$ correction.  Thus no
algorithm that solves all such local-search instances can improve this degree
exponent.  This is evidence about the limits of generic methods, not a claim that
every abstract \LOP\ individually satisfies this lower bound.

\subsection{Weighted normalization and the upper bound}

For each vertex define
\begin{equation}\label{eq:cost}
 b(u)=\max_{z\in\N_s[u]}|\N_s[z]|,
 \qquad c(u)=\frac1{b(u)},
 \qquad c(S)=\sum_{u\in S}c(u).
\end{equation}
These values are computable from radius-$2s$ neighborhoods.

\begin{lemma}[A permitted correction has unit cost]\label{lem:unit-cost}
If $S\subseteq\N_s[v]$, then $c(S)\leq1$.
\end{lemma}

\begin{proof}
For $u\in\N_s[v]$, undirectedness gives $v\in\N_s[u]$, and hence
$b(u)\geq|\N_s[v]|$.  Therefore
$c(S)\leq |S|/|\N_s[v]|\leq1$.
\end{proof}

For a nonempty flip $S$, define its weighted improving ratio as its cut improvement
divided by $c(S)$.  A positive move is weighted minimal if no submove has larger
weighted ratio.  Consider a sequence of weighted minimal moves of ratio at least
$\beta$.  If $w(u)$ counts occurrences, write
\[
 \mu(u)=c(u)w(u),
 \qquad W(U)=\sum_{u\in U}\mu(u).
\]
Since the cut potential ranges over an interval of length at most $|E|$,
\begin{equation}\label{eq:weighted-total}
 \beta W(V)\leq |E|.
\end{equation}
If $v$ occurs, then $W(\{v\})\geq c(v)\geq1/n$.

\begin{lemma}[Weighted replay]\label{lem:weighted-replay}
Replaying the sequence inside $U$ decreases the uncut-edge potential by at least
\[
 \beta W(U)-2\sum_{u\in U}\mu(u)b(u)\deg_{V\setminus U}(u).
\]
\end{lemma}

\begin{proof}
At one step, split the full support into $A\subseteq U$ and $B\subseteq V\setminus
U$.  If the full weighted ratio is $q$, weighted minimality bounds the improvement
of $B$ alone by $qc(B)$.  Hence applying $A$ after $B$ gains at least $qc(A)$.
Freezing the exterior changes the marginal only on crossing edges, with discrepancy
at most two per edge.  Sum over steps and use
$w(u)=\mu(u)/c(u)=\mu(u)b(u)$.
\end{proof}

\begin{lemma}[Weighted sequence localization]\label{lem:weighted-localization}
Let $0<\eps<\beta$ and suppose $v$ occurs in a weighted minimal
$\beta$-improving sequence.  In the initial cut there is a weighted minimal move of
ratio at least $\beta-\eps$, supported in the sequence union and within distance
\begin{equation}\label{eq:weighted-radius}
 O_s\!\left(1+\min\left\{
 \frac{\Delta(\Delta^s+1)}{\eps}\log\!\left(1+\frac{n^3}{\beta}\right),
 \sqrt{\frac{\Delta n}{\eps}\log\!\left(1+\frac{n^3}{\beta}\right)},
 \left(\frac{n^2}{\eps}\log\!\left(1+\frac{n^3}{\beta}\right)\right)^{1/3}
 \right\}\right)
\end{equation}
of $v$.
\end{lemma}

\begin{proof}
Set $h=2s+1$, $U_i=\N_{hi}[v]$, $W_i=W(U_i)$,
$a_i=W_i-W_{i-1}$, and $Z_i=|U_{i+1}\setminus U_{i-1}|$.  A boundary vertex
$u\in U_i$ lies in $U_i\setminus U_{i-1}$.  If it has a neighbor outside $U_i$,
then for every $z\in\N_s[u]$, the ball $\N_s[z]$ lies inside
$U_{i+1}\setminus U_{i-1}$.  Consequently
\[
 b(u)\leq\min\{O_s(\Delta^s),Z_i\},
 \qquad
 \deg_{V\setminus U_i}(u)\leq\min\{\Delta,Z_i\}.
\]
Lemma~\ref{lem:weighted-replay} bounds the replay error by
\begin{equation}\label{eq:weighted-error}
 O_s(a_i)\min\{\Delta(\Delta^s+1),\Delta Z_i,Z_i^2\}.
\end{equation}

By \eqref{eq:weighted-total}, $W(V)\leq n^2/(2\beta)$, while
$W_0\geq1/n$.  Hence
\begin{equation}\label{eq:weighted-harmonic}
 \sum_{i=1}^k\frac{a_i}{W_i}
 \leq\log\frac{W_k}{W_0}
 \leq\log\!\left(1+\frac{n^3}{\beta}\right).
\end{equation}
Also $\sum_{i=1}^kZ_i\leq2n$.  If \eqref{eq:weighted-error} exceeded
$\eps W_i$ for every $i\leq k$, its three terms would give, respectively,
\[
 \log\!\left(1+\frac{n^3}{\beta}\right)
 =\Omega_s\!\left(\frac{\eps k}{\Delta(\Delta^s+1)}\right),
\]
\[
 \log\!\left(1+\frac{n^3}{\beta}\right)
 =\Omega\!\left(\frac{\eps}{\Delta}\sum_{i=1}^k\frac1{Z_i}\right)
 =\Omega\!\left(\frac{\eps k^2}{\Delta n}\right),
\]
and, using convexity of $z\mapsto z^{-2}$,
\[
 \log\!\left(1+\frac{n^3}{\beta}\right)
 =\Omega\!\left(\eps\sum_{i=1}^k\frac1{Z_i^2}\right)
 =\Omega\!\left(\frac{\eps k^3}{n^2}\right).
\]
Thus a good annulus occurs within \eqref{eq:weighted-radius}.  Replaying there gains
at least $(\beta-\eps)W_i$.  Its final support $A$ satisfies $c(A)\leq W_i$, so the
replay has weighted ratio at least $\beta-\eps$; taking a weighted minimal submove
finishes the proof.
\end{proof}

To obtain \eqref{eq:radius-upper}, use the scheduler of Section~4 with
$\eta=1/(8Q)$ and $R_j=1/4+(j-1)/(4Q)$.  A violated radius-$s$ constraint has an
integer cut improvement at least one and, by Lemma~\ref{lem:unit-cost}, weighted
ratio at least one.  The write region for a cluster is enlarged from $\N_L[C]$ to
$\N_{L+s}[C]$.  Compute the decomposition on
$G^p$ for $p=2(L+s+1)+1$ and gather $\N_{L+s+2}[C]$ at the leader.  Thus distinct
same-color write regions are at distance greater than one, and no edge potential is
affected by two simultaneous cluster simulations.  If a final improving coalition
lies in $\N_s[v]$, choose a vertex of that coalition and apply
Lemma~\ref{lem:weighted-localization}; the localized move lies in the enlarged write
region and contradicts exhaustion of the appropriate color phase.  Substituting
$\eps=\Theta(1/\log n)$ in
\eqref{eq:weighted-radius} and multiplying by the $O(\log^2 n\,\poly(\log\log n))$
decomposition overhead gives \eqref{eq:radius-upper}.

\subsection{The lower bound}

The construction adapts the causal indistinguishability strategy used for $\LOC$ in
Section~5 and Appendix~B of Balliu et al.~\cite{BalliuEtAl2026}.  We give the graph
construction and parameter optimization in full.

\begin{lemma}[Adjustable expanding core]\label{lem:core}
For every fixed $s\geq1$, there is $q_0(s)$ such that, for all integers
$q\geq q_0(s)$ and $q\leq M\leq q^s$, there is a simple bipartite graph $C$ with
parts of size $M$, maximum degree at most $2q$, and
\[
 |\partial_C X|\geq \frac q4|X|
 \qquad\text{for every }0<|X|\leq M.
\]
Here $\partial_C X$ is the set of core edges with exactly one endpoint in $X$.
\end{lemma}

\begin{proof}
Include each possible bipartite edge independently with probability $q/M$.  Every
degree has expectation $q$, so a Chernoff bound and a union bound give
$\Pr[\Delta(C)>2q]\leq2M e^{-q/3}$.  If $X$ contains $a$ vertices of one side and
$b$ of the other, with $t=a+b\leq M$, then
\[
 \mathbb E|\partial_CX|
 =\frac qM\bigl(a(M-b)+b(M-a)\bigr)
 =q\left(t-\frac{2ab}{M}\right)
 \geq\frac{qt}{2}.
\]
Another Chernoff bound gives
$\Pr[|\partial_CX|<qt/4]\leq e^{-qt/16}$.  There are at most
$\binom{2M}{t}\leq(2eM/t)^t$ choices of size $t$.  Since
$\log M\leq s\log q=o(q)$ for fixed $s$, the sum of all failure probabilities is
less than one for sufficiently large $q$.
\end{proof}

Turn this core into a radius-$s$ gadget.  Index the vertices on each side by $M$
distinct words from $[q]^s$ and attach the trie of these words to a new root, with
the bipartition side included in the first branch.  The root has degree at most
$2q$, other internal trie vertices have degree at most $q+1$, every core vertex is
at distance exactly $s$ from the root, and the gadget has $\Theta_s(M)$ vertices.
Let $a=\lfloor q/16\rfloor$.  Interfaces will use only side-$0$ core vertices and
will give each such vertex at most $2a+1$ interface edges.  Including its tree-parent
edge, a core vertex has fewer than $q/4$ neighbors outside the core when $q$ is large.

\begin{lemma}[Rigidity]\label{lem:rigidity}
In every radius-$s$ locally optimal cut, a gadget core is labeled by one of its two
proper bipartite colorings.
\end{lemma}

\begin{proof}
Let $\sigma(u)\in\{0,1\}$ denote the bipartition side of a core vertex and call
$x(u)\oplus\sigma(u)$ its phase.  If phases are not constant, let $X$ be the smaller
phase class, so $0<|X|\leq M$.  Every edge of $\partial_CX$ is currently uncut and
becomes cut when $X$ is flipped.  At most $(2a+2)|X|$ non-core edges can be lost.
By Lemma~\ref{lem:core},
\[
 |\partial_CX|\geq(q/4)|X|>(2a+2)|X|
\]
for large $q$.  The move is permitted because the entire core lies in the
radius-$s$ ball of the root, contradicting local optimality.
\end{proof}

The common label on side-$0$ core vertices is the gadget phase.  Flipping every
vertex of a gadget is also a permitted radius-$s$ move.  It changes no internal edge
and toggles precisely the interface edges.

Number the side-$0$ vertices modulo $M$.  For every shift
$t\in\{0,\ldots,a-1\}$, the edges $j\leftrightarrow j+t\pmod M$ form a perfect
matching between two gadgets.  The $a$ matchings are edge-disjoint.  Hence every
integer interface weight at most $aM$ is realized by taking a prefix of these edges,
with interface degree at most $a$ per port vertex.

Let $k=\lfloor aM/2\rfloor$.  Take gadgets
\[
 A_0,\ldots,A_{k-1},
 \qquad B_0,\ldots,B_{k-1},
 \qquad Z.
\]
For $1\leq i\leq k-1$, join $A_{i-1}$ to $A_i$ and $B_{i-1}$ to $B_i$ by
interfaces of weight $2i$.  Form two inputs:
\begin{itemize}
  \item in $G_{\mathrm{odd}}$, join $A_{k-1}$ to $B_{k-1}$ by an interface of
  weight $2k$ and add one edge from $A_{k-1}$ to $Z$;
  \item in $G_{\mathrm{even}}$, join each of $A_{k-1}$ and $B_{k-1}$ to $Z$ by
  an interface of weight $2k$.
\end{itemize}
The graphs are connected, simple, and have maximum degree at most $3q$ for large
$q$.

\begin{lemma}[Forced parity]\label{lem:parity}
Let $u$ and $v$ be side-$0$ representatives in $A_0$ and $B_0$, respectively.
Every radius-$s$ locally optimal cut satisfies
\[
 x(u)\ne x(v)\quad\text{in }G_{\mathrm{odd}},
 \qquad
 x(u)=x(v)\quad\text{in }G_{\mathrm{even}}.
\]
\end{lemma}

\begin{proof}
By Lemma~\ref{lem:rigidity}, every gadget has a phase.  On either arm, an inward
interface has weight two larger than the preceding outward interface.  If the
inward interface were uncut, flipping the whole gadget would gain its full weight
and lose at most the outward weight, a strict improvement.  At the exceptional tip
of $G_{\mathrm{odd}}$, the one extra edge can reduce the gain by only one.  Hence
all main-chain interfaces are cut.

The path of interfaces from $A_0$ to $B_0$ has $2k-1$ links in
$G_{\mathrm{odd}}$ and $2k$ links in $G_{\mathrm{even}}$.  Side-$0$ labels equal
gadget phases, and a cut interface joins opposite phases.  The stated parity follows.
\end{proof}

Give corresponding vertices the same identifiers in the two inputs.  The graphs
differ only at the inner tips and at $Z$.  Any path from $u$ or $v$ to a differing
edge crosses $\Omega_s(k)$ gadget interfaces.  Therefore, for
$T\leq c_s k$, the joint radius-$T$ input neighborhoods of $(u,v)$ are identical in
the two graphs.  The standard causal indistinguishability argument applies to
randomized \LOCAL, and also to quantum-\LOCAL\ with prior entanglement
independent of the input topology; see
\cite{GavoilleKosowskiMarkiewicz2009,AkbariEtAl2025,BalliuEtAl2026}.  The joint
output distribution at $(u,v)$ is the same on both inputs.  If it assigns probability
$p$ to equality, global success is at most $p$ on $G_{\mathrm{even}}$ and at most
$1-p$ on $G_{\mathrm{odd}}$.  One input therefore has success probability at most
$1/2$.  Consequently
\begin{equation}\label{eq:k-lower}
 T=\Omega_s(k)=\Omega_s(qM).
\end{equation}

There are $\Theta(qM)$ gadgets, each with $\Theta_s(M)$ vertices, so the construction
uses $\Theta_s(qM^2)$ vertices.  It may be padded to any larger target size by a path
attached at $Z$, identically in the two inputs.  Choose a sufficiently small constant
multiple
\[
 q=\Theta_s(\min\{\Delta,n^{1/3}\}),
 \qquad
 M=\min\left\{q^s,\left\lfloor\sqrt{\frac{n}{C_s q}}\right\rfloor\right\}.
\]
For sufficiently large feasible parameters, $q\leq M\leq q^s$, the construction
fits in $n$ vertices, and its degree is at most $\Delta$.  Equation
\eqref{eq:k-lower} becomes
\[
 T=\Omega_s\!\left(
 \min\{\Delta^{s+1},\sqrt{\Delta n},n^{2/3}\}
 \right),
\]
which proves \eqref{eq:radius-lower} and completes
Theorem~\ref{thm:radius-cut}.

\section{Discussion}

The direct replay argument isolates the source of the speedup: a sequence of local
improvements should be localized as a sequence, in the original graph, rather than
as a collection of vertices in a second graph.  Minimality controls the discarded
portion of every move.  Once that cost is paid, only an annular boundary remains,
and the two ways of counting that boundary produce the degree and square-root
branches of the running time.

For ordinary locally optimal cut, edge boundaries make the resulting algorithm
polynomially optimal in both $\Delta$ and $n$.  For general local potentials,
Theorem~\ref{thm:generic} is an upper bound rather than a universal characterization:
special algebraic structure may permit faster algorithms.  The radius-local cut
construction shows, however, that larger correction neighborhoods genuinely realize
the boundary-volume exponents and that no generic algorithm covering the entire
class can avoid them.

Three questions remain open.  First, the logarithmic factors in
Corollary~\ref{cor:loc} are not known to be tight.  Second, at constant degree there
is still a gap between deterministic, randomized, and quantum complexities.  Third,
it would be useful to identify structural conditions on a potential---beyond being
an edge sum---that replace the generic local scale $K$ by a smaller intrinsic
boundary parameter.

\begin{thebibliography}{10}

\bibitem{AkbariEtAl2025}
A.~Akbari, X.~Coiteux-Roy, F.~D'Amore, F.~Le Gall, H.~Lievonen,
D.~Melnyk, A.~Modanese, S.~Pai, M.~Renou, V.~Rozho\v{n}, and J.~Suomela.
\newblock Online locality meets distributed quantum computing.
\newblock In \emph{Proceedings of the 57th Annual ACM Symposium on Theory of
Computing (STOC)}, pages 1295--1306, 2025.

\bibitem{BalliuEtAl2026}
A.~Balliu, T.~Boudier, F.~d'Amore, F.~Kuhn, D.~Olivetti, G.~Schmid, and
J.~Suomela.
\newblock Distributed algorithms for potential problems.
\newblock \emph{arXiv:2507.12038v4}, 2026.

\bibitem{GavoilleKosowskiMarkiewicz2009}
C.~Gavoille, A.~Kosowski, and M.~Markiewicz.
\newblock What can be observed locally?
\newblock In \emph{Proceedings of the 23rd International Symposium on Distributed
Computing (DISC)}, volume 5805 of \emph{LNCS}, pages 243--257, 2009.

\bibitem{GhaffariGrunau2024}
M.~Ghaffari and C.~Grunau.
\newblock Near-optimal deterministic network decomposition and ruling set, and
improved {MIS}.
\newblock In \emph{Proceedings of the 65th IEEE Symposium on Foundations of
Computer Science (FOCS)}, pages 2148--2179, 2024.
\newblock Also available as arXiv:2410.19516.

\bibitem{Linial1992}
N.~Linial.
\newblock Locality in distributed graph algorithms.
\newblock \emph{SIAM Journal on Computing}, 21(1):193--201, 1992.

\bibitem{NaorStockmeyer1995}
M.~Naor and L.~J. Stockmeyer.
\newblock What can be computed locally?
\newblock \emph{SIAM Journal on Computing}, 24(6):1259--1277, 1995.

\bibitem{RozhonGhaffari2020}
V.~Rozho\v{n} and M.~Ghaffari.
\newblock Polylogarithmic-time deterministic network decomposition and distributed
derandomization.
\newblock In \emph{Proceedings of the 52nd Annual ACM Symposium on Theory of
Computing (STOC)}, pages 350--363, 2020.

\end{thebibliography}

\end{document}
